\chapter{経緯}\label{ch:history}

作業は 2026 年 10 月 9 日に始まり、21 日までの 13 日間で、リポジトリに 102 のコミットを積んだ（リポジトリには、これとは別に元論文の公式実装の 37 のコミットがある）。この章では、何をどの順でしたか、どこで方針を変えたかを日付の順に書く。各日の最後の括弧は、主なコミットの番号である。

\section{全体の流れ}

\begin{figure}[htbp]
\centering
\begin{tikzpicture}[x=11mm,y=7mm]
\foreach \d/\l in {0/9,1/10,2/11,3/12,4/13,5/14,6/15,7/16,8/17,9/18,10/19,11/20,12/21}{
  \draw[color=ink2!50] (\d,0) -- (\d,-0.15);
  \node[font=\scriptsize,color=ink2] at (\d,-0.45) {10/\l};
}
\draw[color=ink2] (0,0) -- (12,0);
\fill[c1!25] (0,0.3) rectangle (4.9,0.9);  \node[font=\scriptsize] at (2.45,0.6) {再現・データ基盤・実験 A〜G};
\fill[c4!30] (5,0.3) rectangle (6.9,0.9);  \node[font=\scriptsize] at (5.95,0.6) {解説・DB 拡張};
\fill[c7!25] (7,0.3) rectangle (8.3,0.9);  \node[font=\scriptsize] at (7.65,0.6) {H01〜H07};
\fill[c3!30] (8,1.1) rectangle (10.9,1.7);  \node[font=\scriptsize] at (9.45,1.4) {決算発表：実験 2〜5・10};
\fill[c2!25] (8.4,1.9) rectangle (10.9,2.5);  \node[font=\scriptsize] at (9.65,2.2) {BLM：X01〜X05・実験 7〜9};
\fill[c7!25] (10.5,2.7) rectangle (12,3.3);  \node[font=\scriptsize] at (11.25,3.0) {実験 11};
\end{tikzpicture}
\caption{作業の時期。}
\label{fig:hist-gantt}
\end{figure}

\begin{longtable}{P{2.0cm}P{12.6cm}}
\toprule
\textbf{日付} & \textbf{主な出来事} \\
\midrule
\endhead
\bottomrule
\endfoot
10/9（木） & 公式実装を取り込み、OpenAI の API で動くように直した。\code{run.py} を \code{blm} パッケージに分けた（\code{c5fab06}〜\code{2f669bd}） \\
10/10（金） & 設定を \code{.env} から読む形にし、SQLite のデータベースを作った。2023-03〜2026-03 の \SP{} の構成銘柄と株価を入れ、BLM の比率の計算（\code{blm/}）を加えた（\code{6a2dff9}〜\code{e488a76}） \\
10/11〜13 & 実行 1（14,688 件、44 万個の回答）を集め、元論文の計算を再現した。計算の問題点を 9 つ確かめ、データ基盤を作り直した。月次レポート、セクターと銘柄の指標を加え、実験 A〜G を行った。\textbf{総合報告書 1} を書いた（\code{bdd087f}、\code{ff8516b}） \\
10/14（火） & ローカルの LLM の実行の仕組みを加えた。解説「平均分散から Black--Litterman へ」と「BLM のビューとリスクのコツ」を書いた（\code{1a5a0ba}、\code{e7fc413}） \\
10/15（水） & データベースを 2015 年までさかのぼって広げた（発行済株式数は yfinance、足りない分は EDGAR から）（\code{d0d5cf3}） \\
10/16（木） & 仮説 H01〜H07 の設計書を書き、H01〜H03 を実行した。H01 で、元の最適化（SLSQP）が再現できないことを見つけた（\code{95e184d}〜\code{887c451}） \\
10/17（金） & H04〜H07 を実行し、\textbf{実験結果レポート}を書いた。別のモデル（GPT・Gemini）にレビューさせ、指摘を確かめた。ここで\textbf{方針を変えた}：月次のリターンを当てるのをやめ、決算発表への反応を調べる。設計（\code{設計.md}）と、決算発表の取得・時刻の監査・点数付けの仕組みを作った。共分散のノイズ除去（X01）を調べた。Claude・GPT・Gemini に方針と文献調査を独立に書かせて比べた（\code{3f661b3}〜\code{d95db24}） \\
10/18（土） & この研究で最も多い 31 のコミットの日。\textbf{実験 2}（事前登録 → 1 回の実行、+120bp）、アナリスト予想を取得して\textbf{実験 3}（+116bp）、\textbf{実験 4}（FinBERT の追加学習、+102bp）。\textbf{総合報告書 2}。資料集を番号つきのフォルダに整理した。X02〜X04（共分散の窓、TE の校正、VIX）、実験 6 の事前登録の案、EODHD の調査、LLM のシナリオと CVaR の調査、\textbf{実験 7・8}（月次・日次の BLM、支持されない）、\textbf{実験 5}（ChronoGPT の埋め込み、+34bp）（\code{2df746e}〜\code{66ad57f}） \\
10/19（日） & \textbf{実験 9}（約 10 年の月次の BLM、+0.003\%）。日付の誤りを見つけ、直した実行も報告した。\textbf{実験 10}（ChronoGPT の LoRA、7.6 時間の学習、FinBERT を超えない）。撤回された論文（Kim・Muhn・Nikolaev）の引用を直した。X05（$\alpha$ の大きさ）。追加のデータ（提出書類の索引、出来高、ファクター、マクロ、空売り、閲覧数、内部者取引）を取得した。実験 A・B の案を議論し、読ませ方とまとめ方を Deep Research で調べた（\code{0856fd0}〜\code{ba198ac}） \\
10/20（月） & 修正 1〜5 の方針、事前登録、コード。修正 1 の規則の穴（行ごと捨てていた）を、評価の前に見つけて直した。RunPod の GPU で点数付けを始めたが、結果をダウンロードする前に Pod を止めて消えた（\code{71a5fec}〜\code{ac0f458}） \\
10/21（火） & この PC で約 19 時間かけて点数を付け直し、開発の判定と最後の比較を 1 回ずつ行った。\textbf{実験 11} のレポートを書いた。この本をまとめた（\code{ff35527}、\code{6591f06}） \\
\end{longtable}

\section{方針の分かれ目}

\subsection{10 月 13 日：数字をなぞる LLM}

総合報告書 1 の時点で、gpt-4o-mini は、どう聞いても渡した数字の中の直近の値動きに沿って答えることがわかった。7 つの実験（A〜G）で、入力を要約する、財務を足す、比率を明言する、会話にする、5 段階で答えさせる、銘柄の数字を渡さない、を試したが変わらなかった。結論は「数字の処理は統計モデルに任せ、LLM には文章の解釈だけをさせる」だった。

\subsection{10 月 17 日：月次のリターンから、決算発表の反応へ}

H01〜H07 で、元論文の形の弱さを統計的に確かめた（H02：LLM に固有の予測力はあっても IC 0.03 未満）。H06 で、決算発表文の点数は値動きをなぞらないが、翌月のリターンは説明しなかった。そこで問いを変えた。「LLM の点数で翌月を当てる」のではなく、「LLM の点数は、発表への株価の反応を説明するか」を先に確かめる。反応を説明しないなら、翌月を当てられるはずがない。反応を説明するなら、次にその情報がいつまで残るかを調べればよい。

\subsection{10 月 18 日：事前登録}

実験 2 から、\textbf{株価の反応を一度も計算する前に}、計画と評価のコードを同じコミットで固定し、1 回だけ実行する方式にした。これで「結果を見てから都合のよい分析を選ぶ」ことができなくなる。実験 2〜11 と X02〜X04 はすべてこの方式で行った（X05 は結果を見た後の探索で、そう明記した）。

\subsection{10 月 18 日：前向きの検証をやめる}

総合報告書 2 の計画（GPT・Gemini の立案を比べたもの）では、結果を見ていない 2026 年 10〜11 月の決算発表で前向きに再現することを最優先にしていた。しかし 12 月まで待つ必要があるので、過去のデータでの検証だけで進めることにした（利用者の判断）。

\subsection{10 月 19 日：深層学習をどこまで効かせられるか}

実験 4 で、反応で追加学習した FinBERT は +71bp と強かった。ChronoGPT も中まで学習させれば伸びるかを、LoRA で確かめた（実験 10）。結果は、線形のリッジ（実験 5）より弱かった。

\subsection{10 月 19〜20 日：スコープを絞る}

追加のデータを集めたあと、「横に広げすぎていないか」を見直した。スコープを「LLM の読み取りで何が効いたか」と「読ませ方とまとめ方で直せる場所はあるか」に絞り、gpt-4o-mini は最後の比較にだけ使い、試すのは ChronoGPT と従来の NLP にする、と決めた。これが実験 11 になった。

\section{コミットの数}

\begin{table}[htbp]
\centering
\caption{日ごとのコミットの数（この研究の分だけ）。}
\small
\begin{tabular}{lrrrrrrrrrr}
\toprule
日付 & 10/9 & 10/10 & 10/13 & 10/14 & 10/15 & 10/16 & 10/17 & 10/18 & 10/19 & 10/20--21 \\
\midrule
コミット & 4 & 9 & 2 & 4 & 2 & 4 & 13 & 31 & 20 & 13 \\
\bottomrule
\end{tabular}
\end{table}

10 月 11・12 日は、データの取得と実行 1 の回答の収集に時間を使い、コミットは 13 日にまとめた。
